% Framework section removed as it is covered in Part I.

\section{Synthesis of Constructive Estimates}
\label{sec:technical_details_app145}

\subsection{Strong Coupling Cluster Expansion}
In the regime of small $\beta$ (strong coupling), we utilize the character expansion of the Boltzmann weight. For $U \in SU(N)$:
\begin{equation}
e^{\frac{\beta}{N} \mathrm{Re}\,\mathrm{Tr}(U)} = c_0(\beta) \left( 1 + \sum_{r \neq 0} d_r a_r(\beta) \chi_r(U) \right)
\end{equation}
where the sum runs over non-trivial irreducible representations $r$, $d_r$ is the dimension, and $a_r(\beta) = \frac{u_r(\beta)}{u_0(\beta)}$ are the expansion coefficients involving modified Bessel functions (for $U(1)$) or their non-abelian generalizations.

\begin{definition}[Polymer System]
A polymer $\gamma$ is a connected set of plaquettes. Two polymers $\gamma_1, \gamma_2$ are incompatible if they share a link. The partition function admits the representation:
\begin{equation}
Z_\Lambda = \sum_{\Gamma} \prod_{\gamma \in \Gamma} W(\gamma)
\end{equation}
where the sum is over sets $\Gamma$ of mutually compatible polymers.
\end{definition}

\begin{theorem}[Convergence of Cluster Expansion]
\label{thm:cluster_convergence_app145}
\label{thm:cluster-expansion}
There exists $\beta_0 > 0$ such that for all $\beta < \beta_0$, the polymer weights satisfy the Koteck\'y-Preiss condition:
\begin{equation}
\sum_{\gamma \ni p} |W(\gamma)| e^{a(\beta) |\gamma|} \le 1
\end{equation}
for some $a(\beta) > 0$. Consequently, the free energy density $f(\beta) = \lim_{\Lambda \to \infty} |\Lambda|^{-1} \ln Z_\Lambda$ is analytic in $\beta$.
\end{theorem}

\begin{proof}
The weight $W(\gamma)$ is obtained by integrating the product of character expansion coefficients over the links in $\gamma$. For a polymer to have non-zero weight, every link must belong to at least two plaquettes (or zero, but polymers are connected sets of plaquettes) such that the tensor product of representations contains the trivial representation.
For small $\beta$, $a_r(\beta) = O(\beta^{k_r})$. The leading order contribution comes from the fundamental representation.
We have the bound $|W(\gamma)| \le (C \beta)^{|\gamma|}$.
The number of polymers of size $n$ containing a fixed plaquette $p$ is bounded by $(2d e)^n$.
Thus, $\sum_{\gamma \ni p} |W(\gamma)| e^{|\gamma|} \le \sum_{n \ge 6} (2d e)^n (C \beta)^n e^n$.
For $\beta < (2de^2 C)^{-1}$, this series converges and is small, satisfying the condition.
\end{proof}

\subsection{Mass Gap in Strong Coupling}
\begin{theorem}[Exponential Decay of Correlations]
For $\beta < \beta_0$, the two-point correlation function of the plaquette variables decays exponentially:
\begin{equation}
|\langle \chi(U_p) \chi(U_{p'}) \rangle - \langle \chi(U_p) \rangle \langle \chi(U_{p'}) \rangle| \le K e^{-m(\beta) \|p - p'\|}
\end{equation}
where $m(\beta) = - \ln(C \beta) + O(\beta) > 0$ is the mass gap.
\end{theorem}

\begin{proof}
The truncated correlation function is given by the sum over all polymers $\gamma$ connecting $p$ and $p'$ in the cluster expansion of the observable insertion.
Let $C_{p,p'}$ be the class of polymers connecting $p$ and $p'$.
\begin{equation}
\langle \mathcal{O}_p \mathcal{O}_{p'} \rangle_c = \sum_{\gamma \in C_{p,p'}} W(\gamma) \Xi(\gamma)
\end{equation}
where $\Xi(\gamma)$ represents the modification of the background partition function.
Since $|W(\gamma)| \le (C \beta)^{|\gamma|}$ and $|\gamma| \ge \|p - p'\|$, we have:
\begin{equation}
|\langle \mathcal{O}_p \mathcal{O}_{p'} \rangle_c| \le \sum_{L \ge \|p-p'\|} N(L) (C \beta)^L \le \sum_{L} \mu^L (C \beta)^L
\end{equation}
Convergence requires $C \beta \mu < 1$. The decay rate is governed by the smallest $L$, yielding exponential decay with mass $m \approx -\ln(C \beta)$.
\end{proof}

\subsection{Intermediate Coupling and RP Monotonicity}
To extend the result beyond the radius of convergence of the strong coupling expansion, we employ Reflection Positivity Monotonicity and the Adjoint QCD interpolation.

\begin{theorem}[Unconditional RP Monotonicity]
The string tension $\sigma(\beta)$ is a monotonically decreasing function of $\beta$.
Furthermore, $\sigma(\beta) > 0$ for all $\beta < \infty$.
\end{theorem}

\begin{proof}
See Theorem~\ref{thm:rp-monotonicity-vortex} for the derivation of monotonicity.
The strict positivity $\sigma(\beta) > 0$ is established unconditionally by the Adjoint QCD interpolation (Theorem~\ref{thm:adjoint-gap-main} in Appendix~\ref{app:adjoint-proof}).
The interpolation connects the theory to a regime where Center Symmetry is explicitly broken by the adjoint mass term, preventing the string tension from vanishing, and then recovers the pure gauge theory in the decoupling limit while preserving the gap.
\end{proof}

\subsection{Continuum Limit and Asymptotic Freedom}
The continuum limit is defined by taking $g \to 0$ ($\beta \to \infty$) while adjusting the lattice spacing $a(g)$ such that the physical mass $m_{phys}$ remains constant.
The renormalization group equation for the coupling is:
\begin{equation}
a \frac{dg}{da} = -\beta_0 g^3 - \beta_1 g^5 + O(g^7)
\end{equation}
with $\beta_0 = \frac{11}{3} \frac{N}{16\pi^2} > 0$.
Integrating this yields the scaling law:
\begin{equation}
m_{phys} a \sim \exp\left( -\frac{1}{2\beta_0 g^2} \right)
\end{equation}
This implies that as $g \to 0$, the correlation length $\xi = 1/(m_{phys} a)$ diverges exponentially in lattice units.

\begin{theorem}[Stability of the Mass Gap]
Under the renormalization group flow, the effective action remains in the domain of attraction of the trivial fixed point (high temperature sink) for the effective variables, but the correlation length grows.
Specifically, for sufficiently large $\beta$, the theory is perturbatively close to the free theory but with non-perturbative corrections that generate a mass gap.
The mass gap $m > 0$ persists in the continuum limit $a \to 0$.
\end{theorem}

\begin{proof}
We employ the Balaban block-spin transformation. The effective action $S_{eff}$ after $k$ steps is defined on a lattice of spacing $L^k a$.
The flow of the coupling constant $g_k$ is driven by the beta function.
Since $\beta_0 > 0$, the running coupling $g_k$ grows as we move to the infrared (larger scales).
Eventually, $g_k$ becomes large enough to enter the domain of the strong coupling expansion (Section 2.1).
Once in the strong coupling regime, the mass gap is established by Theorem 2.2.
Since the RG transformation is analytic and preserves the spectral properties (up to rescaling), the existence of the gap at the effective scale implies the existence of the gap at the microscopic scale.
\end{proof}

% Appendices removed.
